\documentclass[a4paper,twoside,12pt]{report}
% Seboholi Mamello (2022)

% Packages
\usepackage{fullpage}
\usepackage{float}
\usepackage{url}
\usepackage{charter}
\usepackage{graphicx}
\usepackage{subfigure}
\usepackage{amsmath}
\usepackage{amssymb}
\usepackage{amsthm}
\usepackage{booktabs}
\usepackage{parskip}
\usepackage[all]{nowidow}
\setnoclub[2]
\setnowidow[2]

\usepackage{dirtytalk}
\usepackage{pdfpages}

% Referencing
\usepackage[]{hyperref}
\usepackage[nameinlink]{cleveref}
\hypersetup{
  colorlinks = true,
  urlcolor = blue,
  linkcolor = blue,
  citecolor = blue
}

% University of the Witwatersrand Citation Style
\usepackage{natbib} % Force natbib.sty to put citation labels in the reference list
\makeatletter
\renewcommand\NAT@biblabel[1]{\def\citeauthoryear##1##2{##1 ##2}[#1]\hfill}
\renewcommand\NAT@bibsetup[1]{%
  \setlength{\itemsep}{\bibsep}\setlength{\parsep}{\z@}}
\def\@lbibitem[#1]#2{%
  \if\relax\@extra@b@citeb\relax\else
    \@ifundefined{br@#2\@extra@b@citeb}{}{%
     \@namedef{br@#2}{\@nameuse{br@#2\@extra@b@citeb}}}\fi
   \@ifundefined{b@#2\@extra@b@citeb}{\def\NAT@num{}}{\NAT@parse{#2}}%
   \item[\hfil\hyper@natanchorstart{#2\@extra@b@citeb}\@biblabel{#1}%
    \hyper@natanchorend]%
    \NAT@ifcmd#1(@)(@)\@nil{#2}}
\makeatother


\bibliographystyle{named-wits}
\bibpunct{[}{]}{;}{a}{}{}  % to get correct punctuation for bibliography
\setlength{\skip\footins}{1.5cm}
\newcommand{\citets}[1]{\citeauthor{#1}'s \citeyearpar{#1}}
\renewcommand\bibname{References} 

\pagestyle{headings}

\renewenvironment{abstract}{\ \vfill\begin{center}\textbf{Abstract}\end{center}\addcontentsline{toc}{section}{Abstract}}{\vfill\vfill\newpage}
\newenvironment{declaration}{\ \vfill\begin{center}\textbf{Declaration}\end{center}\addcontentsline{toc}{section}{Declaration}}{\vfill\vfill\newpage}
\newenvironment{acknowledgements}{\ \vfill\begin{center}\textbf{Acknowledgements}\end{center}\addcontentsline{toc}{section}{Acknowledgements}}{\vfill\vfill\newpage}

\pagestyle{plain}

% Document
\begin{document}
\onecolumn
\thispagestyle{empty}

\pagenumbering{roman}
\setcounter{page}{0}
\addcontentsline{toc}{chapter}{Preface}
\
\begin{center}
  \vfill
  \huge \textbf{Goal Selection Strategies for Learning Goal-Oriented Value Functions}\\[5pt]
  \LARGE \textbf{Reinforcement Learning}\\[20pt]
  \large School of Computer Science \& Applied Mathematics\\
  \large University of the Witwatersrand\\[20pt]
  \large Mamello Seboholi\\
  \large \textbf{1851317}\\[20pt]
  \large Supervised by Dr. Steven James and Prof. Benjamin Rosman\\[10pt]
  \today
  \vfill
  \includegraphics[width=5cm]{images/wits-logo.png}
  \vfill
  \normalsize A proposal submitted to the Faculty of Science, University of the Witwatersrand, Johannesburg, for the 
  Introduction to Research Methods course.
  \vfill
\end{center}
\newpage

\pagestyle{plain}
\setcounter{page}{1}

\phantomsection
\begin{abstract}
Current state of the art in reinforcement learning show impressive outcomes in high-dimensional environments, agents 
are shown to compose new tasks by combining previously learned tasks using boolean algebra. We propose integrating 
Thompson sampling and Upper Confidence Bounds (UCB1) to the Q-value functions as well as the extended Q-value functions, 
to balance between exploration and exploitation.
\end{abstract}

\phantomsection
\addcontentsline{toc}{section}{Declaration}
\includepdf[pages=-]{declaration-form.pdf}

\phantomsection
\addcontentsline{toc}{section}{Table of Contents}
\tableofcontents
\newpage


\pagenumbering{arabic}


\chapter{Introduction}
Reinforcement learning (RL) has as of late seen advances in complex, multi-dimensional problems 
\citep{mnih2015human,levine2016end,lillicrap2015continuous,silver2017mastering}. However, this comes with a challenge 
of practicality as these methods need to be trained with a vast amount of samples. One of the solutions to this problem 
is the use of $composition$ \citep*{todorov2009compositionality} to transfer agent knowledge. This allows the agents to 
build out skills from pre-existing skills or to speed up the training of new skills.

\citet*{nangue2020boolean} built on the works of \citet*{haarnoja2018composable}, \citet*{van2019composing}, and 
\citet*{hunt2019composing} by formalizing the union, intersection and negation of tasks. This allows for zero-shot
composition of new tasks. The Q-value function used utilises greedy based algorithms. We are interested in bandit based 
algorithms, particularly in Thompson sampling \citep{thompson1933likelihood} and Upper Confidence Bounds (UCB1) 
\citep{auer2002finite}.

The proposal takes the following structure. In, \Cref{ch:review}, we review the current literature. We then present the 
problem statement, along with posing the research hypothesis, objectives as well as methods and limitations in 
\Cref{ch:researchMethodology}. In \Cref{ch:plan}, we outline the expected timeline and lastly we summarize the 
proposal in \Cref{ch:conclusion}.


\chapter{Background and Related work} \label{ch:review}

\section{Introduction}
This chapter serves to present our findings in related work to this paper. The paper has the following layout: 
\Cref{review:rl} explains the reinforcement learning problem we are focused in. 
\Cref{review:composition} explores knowledge transfer through composition with the focus on concurrent composition.
\Cref{review:exp} discusses the literature on the exploration-exploitation dilemma; and 
\Cref{review:conclusion} will serve as a conclusion for this chapter.

\section{Reinforcement Learning} \label{review:rl}
\subsection{Markov Decision Processes}
We focus on Reinforcement Learning tasks that are modelled using Markov Decision Processes (MDPs) where an MDP is 
defined as a quadruple (4-tuple) composed of the (i) state space $S$, (ii) action space $A$, (iii) Markov kernel defined 
by $\rho$ that takes $S\times A$ to $S$, and (iv) reward function $r$ that has real values and bounded by a minimum and 
maximum $r$ values.

\subsubsection{Policies}
RL agent's goal is to work out a policy $\pi$ that maps $S$ to $A$, which solves a given task optimally. 

\subsubsection{Value Functions}
Extended reward function and extended Q-value functions are defined. Tasks, as well as the extended Q-value functions 
are defined as Boolean algebra.

\section{Composition} \label{review:composition}
It is essential that we discuss the idea of $composition$ \citep{todorov2009compositionality} within the context of 
this project as we aim to improve on the literature focused on this topic.

\citet{nangue2020boolean} discusses the need for reinforcement learning (RL) agents to have an ability of knowledge 
transfer as reinforcement learning (RL) problems become expensive. 

This paper, along with \citet{todorov2009compositionality}, \citet{saxe2017hierarchy}, \citet{haarnoja2018composable}, 
\citet{van2019composing}, \citet{hunt2019composing}, and \citet{peng2019mcp} focus on concurrent composition, where 
novel tasks are formed by combining previously learned tasks. Another group of literature focus on sequentially chaining 
learned policies to solve complex tasks, examples of this is options \citep{sutton1999between} and hierarchical 
reinforcement learning \citep{barto2003recent}.

This paper is the basis for this project with the focus being on how agents choose between maximizing goals and 
discovering new goals in the environment, which is currently done in a greedy fashion.

\section{Exploration-exploitation} \label{review:exp}
The balance between exploration and exploitation is a well known problem in RL. and much of this project is focused on 
idea. It has been researched rigorously with respect to finding the optimal policy for actions, however, there is no 
significant work relating to the balance between exploration and exploitation for goal-oriented RL.

The literature is usually classified into two types of methods. (i) $undirected$ methods, where agents resolve the 
exploration-exploitation dilemma using Q-values, these types of methods seem to perform very well with small to medium 
problem sizes but do not seem to find optimal policies when the problem scales significantly. (ii) $directed$ methods, 
where agents resolve the exploration-exploitation dilemma by using knowledge about exploration, these methods deal well 
with increased scale of problems, however, this comes with considerable computation requirements.

\subsection{Epsilon Greedy}
Epsilon greedy ($\epsilon$-greedy) is a very simple and popular method used to balance between exploration and 
exploitation. It is an example of an $undirected$ method. It explores the environment with a probability of $\epsilon$,
and chooses an action giving the highest reward with a probability of 1 - $\epsilon$, where $\epsilon$ is chosen to be a 
very small number within the open interval $(0, 1)$ (Note: An $\epsilon$ of 0 means exploitation only and an $\epsilon$ 
of 1 means exploration only). It's origins are not clear, however, it has been used as far back as 
\citet*{watkins1989learning} and \citet*{sutton1998introduction}. Papers like \citet*{tokic2011value} and 
\citet*{dos2017adaptive} have extended the $\epsilon$-greedy method with an attempt to control the exploration rate 
($\epsilon$). An issue with $\epsilon$-greedy is that, during exploration, it chooses equally among the other actions, 
this is not good when there exists actions with negative rewards.

\subsection{Upper Confidence Bounds}
Upper Confidence Bounds (UCB) and specifically UCB1 \citep*{auer2002finite} is a stochastic method that is based on 
optimism. In this method, data is gathered and then used to assign a weight to each arm in the multi-armed bandit 
problem, this weight is known as the upper Confidence bound. UCB methods move focus from exploration to exploitation, 
$\log_{}{t} / \operatorname{N}_t(a)$ is used encourage exploration of the environment as $\operatorname{N}_t(a)$ remains
small for actions that have not been explored for a long time, a parameter $c$ is also used to control level of 
exploration.

\subsection{Thompson Sampling}
\citet*{wyatt1998exploration} introduces Q-value sampling which is a $directed$ method also known as a $stochastic$
method where the rewards are represented by a probability distribution. The probability distribution takes into account 
both exploitation (expected reward) and exploration (how uncertain it is for actual reward). An agent then takes an 
action based on this probability distribution. An issue with this method is that it requires vast amount of data to 
build the probability distribution. \citet*{dearden1998bayesian} uses Q-value sampling along with Myopic value of 
imperfection \citet*{howard1966information} to present a Bayesian based Q-learning method where actions also depend on 
probability distribution. Thompson Sampling \citep*{thompson1933likelihood} is another sampling based algorithm.

\section{Conclusion} \label{review:conclusion}
This chapter serves to showcase literature in Reinforcement Learning and how they have an effect on our research project.
These papers will help in determining an appropriate method that will balance between exploration and exploitation in 
the Goal-oriented Q-learning algorithm for the boolean algebra tasks.


\chapter{Research Methodology} \label{ch:researchMethodology}
\section{Problem Statement}

\section{Hypothesis}
The project's aspiration is to extend the goal-oriented value function by integrating non greedy methods for balancing 
exploration and exploitation with a focus on multi-armed bandit algorithms.

We propose that using Thompson sampling and Upper Confidence Bounds (UCB1) decreases the number of samples required to reach 
convergence as compared to epsilon-greedy. We also hypothesise that for the same sample size, the proposed methods have 
a lesser regret than epsilon-greedy method.

\section{Research Questions}
The above propositions raises the following research questions:

\begin{itemize}
  \item Can we utilise Thompson sampling for the extended Q-value functions defined for Boolean algebra?
  \item Does Thompson sampling based extended Q-value functions reach convergence? Do they require a lesser sample size 
  compared to epsilon-greedy?
  \item Is the regret for Thompson sampling over a range of sample sizes smaller than that of epsilon-greedy for the same
  sample size?
  \item Can UCB1 be used to train the extended Q-value functions defined for Boolean algebra?
  \item Does training reach convergence when using UCB1 and does it reach it with a lesser sample size than epsilon-greedy?
  \item Does UCB1 yield lesser regret compared to epsilon-greedy given that the sample size remains the same?
  \item Which of the proposed algorithms better work with goal-oriented value functions?
\end{itemize}

\section{Methodology}
\subsection{Research Design}

\subsection{Methods}

\subsection{Limitations}

\chapter{Research Plan} \label{ch:plan}

\chapter{Conclusion} \label{ch:conclusion}

\nocite{*}

\bibliography{references}


\end{document}
